\documentclass[12pt,reqno]{amsart}
\usepackage{amsmath,amssymb}
\usepackage{mathptmx}
\usepackage[T1]{fontenc}
\usepackage{enumitem}
\pagestyle{plain}

\setlength{\oddsidemargin}{0in}
\setlength{\evensidemargin}{0in}
\setlength{\topmargin}{0in}
\setlength{\headheight}{0pt}
\setlength{\headsep}{0pt}
\setlength{\footskip}{.3in}
\setlength{\textheight}{9.2in}
\setlength{\textwidth}{6.5in}
\setlength{\parindent}{0pt}
\setlist[enumerate,1]{leftmargin=*,itemsep=1.1em}
\setlist[enumerate,2]{label=(\alph*),leftmargin=1.7em,itemsep=.4em,topsep=.3em}

\newcommand{\C}{\mathbb{C}}
\newcommand{\R}{\mathbb{R}}
\newcommand{\Z}{\mathbb{Z}}

\title{Problem Set 6 \\ MATH 411: Honors Complex Variables, Fall 2026}

% Coverage: Lectures 27--30 (analytic isomorphisms, Schwarz's lemma,
% automorphisms of the disc, Fourier transform, and Poisson summation).
% Sources:
% 1--2: Old Homework/Homework/Homework7_Fall2021.tex, Q1
%       (Lang, Complex Analysis, 4th ed., Ch. VII, Section 2, Exercises 1--2).
% 3:    Old Homework/Homework/Homework7_Fall2021.tex, Q3
%       (also Old Homework/Homework/hwk4.tex, Q9).
% 4--5: Stein & Shakarchi, Complex Analysis, Ch. 4, Exercises 1 and 7.

\begin{document}
\maketitle

These are practice problems and are not to be turned in. Justify your
answers and all limiting arguments. Unless stated otherwise, closed
contours are oriented counterclockwise.

Suggested completion date: Friday, December 4, 2026.

\begin{enumerate}

\item \textbf{Fixed points of maps of the unit disc.}
Let $D=\{z\in\C:|z|<1\}$ be the open unit disc, and let $f:D\to D$ be
holomorphic. Prove that if there exist two distinct points $a,b\in D$
that are fixed by $f$---that is, $f(a)=a$ and $f(b)=b$---then $f(z)=z$
for all $z\in D$.
\emph{Hint:} Conjugate $f$ by an analytic automorphism of $D$ sending
$0$ to $a$, and apply Schwarz's lemma.


\item \textbf{The Schwarz--Pick lemma.}
Let $f:D\to D$ be a holomorphic map of the open unit disc into itself.
Prove that for all $a\in D$,
\[
 \frac{|f'(a)|}{1-|f(a)|^2}\leq\frac{1}{1-|a|^2}.
\]
\emph{Hint:} Let $g$ be an analytic automorphism of $D$ with $g(0)=a$,
and let $h$ be an analytic automorphism of $D$ with $h(f(a))=0$. Set
$F=h\circ f\circ g$, compute $F'(0)$, and apply Schwarz's lemma.


\item \textbf{The Joukowski map.}
Consider the map $f(z)=z+1/z$.
\begin{enumerate}
\item Show that $f$ is an analytic isomorphism from the exterior of the
closed unit disc, $U=\{z\in\C:|z|>1\}$, onto the slit plane
$\C\setminus[-2,2]$.
\item What is the image of the unit circle $\{|z|=1\}$ under $f$?
\item In polar coordinates $z=re^{i\theta}$, write $w=f(z)=u+iv$.
Verify that
\[
 u=\left(r+\frac{1}{r}\right)\cos\theta
 \qquad\text{and}\qquad
 v=\left(r-\frac{1}{r}\right)\sin\theta.
\]
Show that for each $c>1$, the circle $r=c$ maps onto an ellipse with
semi-major axis $c+1/c$ and semi-minor axis $c-1/c$. Show also that
each radial ray $\theta=c$ ($r>1$, with $\cos c\sin c\neq0$) maps onto
one quarter of a hyperbola.
\end{enumerate}



\item \textbf{Uniqueness of the Fourier transform.}
Suppose $f:\R\to\C$ is continuous and of moderate decrease (that is,
$|f(x)|\leq A/(1+x^2)$ for some constant $A>0$ and all $x\in\R$), and
suppose that its Fourier transform vanishes identically:
\[
 \widehat f(\xi)=\int_{-\infty}^{\infty}f(x)e^{-2\pi i x\xi}\,dx=0
 \qquad\text{for all }\xi\in\R.
\]
Show that $f=0$ by completing the following outline.
\begin{enumerate}
\item For each fixed real number $t\in\R$, consider the two functions
\[
 A(z)=\int_{-\infty}^{t} f(x)e^{-2\pi i z(x-t)}\,dx
 \qquad\text{and}\qquad
 B(z)=-\int_{t}^{\infty} f(x)e^{-2\pi i z(x-t)}\,dx.
\]
Show that $A(\xi)=B(\xi)$ for all $\xi\in\R$.
\item Prove that the function $F$ equal to $A$ in the closed upper
half-plane $\{\operatorname{Im}z\geq 0\}$ and to $B$ in the lower
half-plane $\{\operatorname{Im}z<0\}$ is entire and bounded, hence
constant. In fact, show that $F=0$.
\item Deduce that $f=0$.
\end{enumerate}


\item \textbf{A summation identity via Poisson summation.}
Let $\tau\in\C$ satisfy $\operatorname{Im}\tau>0$, and let $k\geq 2$ be
an integer.
\begin{enumerate}
\item Consider $f(z)=1/(\tau+z)^k$. Show that $f\in\mathcal{F}_a$
for some $a>0$, compute its Fourier transform $\widehat f(\xi)$ for all
$\xi\in\R$ using contour integration and residues, and apply the Poisson
summation formula to prove that
\[
 \sum_{n=-\infty}^{\infty}\frac{1}{(\tau+n)^k}
 =\frac{(-2\pi i)^k}{(k-1)!}\sum_{n=1}^{\infty} n^{k-1} e^{2\pi i n\tau}.
\]
\item Setting $k=2$, deduce that
\[
 \sum_{n=-\infty}^{\infty}\frac{1}{(\tau+n)^2}
 =\frac{\pi^2}{\sin^2(\pi\tau)}
\]
for all $\tau\in\C$ with $\operatorname{Im}\tau>0$ (and hence, by
analytic continuation, for all $\tau\in\C\setminus\Z$).
\end{enumerate}

\end{enumerate}
\end{document}
